\documentclass[10pt,a4paper]{article}
	
	\usepackage{amsmath,amssymb,amsthm,mathtools,amsfonts}
	\usepackage{breqn}
	\usepackage[utf8]{inputenc}
	\usepackage[english]{babel}
	\usepackage[T1]{fontenc}
	\usepackage{graphicx}
	\graphicspath{{./img/}}
	
	%\usepackage{ebgaramond}
	\usepackage{fontspec}
	\setmainfont[Numbers=OldStyle]{EB Garamond}
	\newfontfamily\plantinb[Numbers=OldStyle]{Plantin MT Pro Bold Condensed}
	\newfontfamily\plantin[Numbers=OldStyle]{Plantin MT Pro Semibold}
	
	\usepackage[pdfborder={0 0 0}]{hyperref} %% Ingen firkant rundt om referencer
	\usepackage{multicol}
	\usepackage{tabularx}
	\usepackage{enumitem}% better controls with enumerating
	\usepackage{wrapfig}
		%\setlength{\wrapoverhang}{\marginparwidth}
		%\addtolength{\wrapoverhang}{\marginparsep}
		\setlength{\columnsep}{0pt}
		\setlength{\intextsep}{-10pt}
	\usepackage{caption}
	\usepackage[svgnames]{xcolor}
	\usepackage{dirtytalk}
	\usepackage[modulo]{lineno}
	\usepackage{hyperref}
	\usepackage{tasks}
	\newcommand\svar[1]{\newline\textcolor{red}{\textbf{Svar}\\#1}}
	\newcommand{\qm}[1]{‘‘#1”}
	\newcommand{\sqm}[1]{‘#1’}
	
	\definecolor{hgblue}{RGB}{10, 40, 80}
	\definecolor{hgred}{RGB}{140, 10, 10}
	\definecolor{hggrey}{RGB}{215, 215, 210}
	\definecolor{hggreen}{RGB}{145, 205, 205}
	
	\usepackage{titlesec}
	\titleformat{\subsection}{\color{hgred}\mysectionfont\large}{}{}{}
\newcommand{\source}[1]{\footnote{Source: \href{#1}{#1}}}
	\usepackage{lipsum}
	\usepackage{comment}
	\usepackage{ulem}
	
	\usepackage{bigfoot}
		\DeclareNewFootnote{a}
			\renewcommand{\thefootnotea}{\fnsymbol{footnotea}}
				\newcommand{\footnotes}[1]{\footnotea[2]{#1}} % here you can specify what symbol you want in footnotes
				\newcommand{\footnoted}[1]{\footnotea[3]{#1}} % for a second footnote on a page
				\MakePerPage{footnotes}
	
	%\reversemarginpar
	\newcommand{\glos}[3]{\dotuline{#1}\marginpar{\scriptsize\textit{#3} #2}}

	\setlength{\parskip}{0.5\baselineskip}
	\setlength{\parindent}{0cm}
	\setlist[enumerate]{label=\alph*),left=20pt}
	%\setlist[enumerate]{leftmargin=\parindent}
	%\setlist[enumerate]{nosep}
	
	%Titling
	\usepackage{titling}
	\setlength{\droptitle}{-12ex}
	\pretitle{\begin{flushleft}\plantinb\huge\color{hgblue}}
	\posttitle{\par\end{flushleft}}
	\preauthor{\begin{flushleft}\Large}
	\postauthor{\end{flushleft}}
	\predate{\begin{flushleft}}
	\postdate{\end{flushleft}}
	
	\setcounter{secnumdepth}{0}
	
	\usepackage{tikz}

\begin{document}
%\pagecolor{FloralWhite}
\title{The Coronation of Mr. Thomas Shap}
%\date{\vspace{-3em}} % I tilfældet ingen dato
\date{1912}
\author{Lord Dunsany}
\maketitle
    
\thispagestyle{plain}
\setcounter{page}{1}
\linenumbers It was the occupation of Mr. Thomas Shap to persuade customers that the goods were genuine and of an excellent quality, and that as regards the price their unspoken will was consulted. And in order to carry on this occupation he went by train very early every morning some few miles nearer to the City from the suburb in which he slept. This was the use to which he put his life.

From the moment when he first perceived (not as one reads a thing in a book, but as truths are revealed to one's instinct) the very beastliness of his occupation, and of the house that he slept in, its shape, make and pretensions, and even of the clothes that he wore; from that moment he withdrew his dreams from it, his fancies, his ambitions, everything in fact except that \glos{ponderable}{physical; able to be weighed or measured}{adj.} Mr. Shap that dressed in a frock-coat, bought tickets and handled money and could in turn be handled by the statistician. The priest's share in Mr. Shap, the share of the poet, never caught the early train to the City at all.

He used to take little flights of fancy at first, dwelt all day in his dreamy way on fields and rivers lying in the sunlight where it strikes the world more brilliantly further South. And then he began to imagine butterflies there; after that, silken people and the temples they built to their gods.

They noticed that he was silent, and even absent at times, but they found no fault with his behaviour with customers, to whom he remained as \glos{plausible}{convincing or believable}{adj.} as of old. So he dreamed for a year, and his fancy gained strength as he dreamed. He still read halfpenny papers in the train, still discussed the passing day's \glos{ephemeral}{short-lived; lasting only a brief time}{adj.} topic, still voted at elections, though he no longer did these things with the whole Shap—his soul was no longer in them.

He had had a pleasant year, his imagination was all new to him still, and it had often discovered beautiful things away where it went, southeast at the edge of the twilight. And he had a matter-of-fact and logical mind, so that he often said, "Why should I pay my twopence at the electric theatre when I can see all sorts of things quite easily without?" Whatever he did was logical before anything else, and those that knew him always spoke of Shap as "a sound, sane, level-headed man."

On far the most important day of his life he went as usual to town by the early train to sell plausible articles to customers, while the spiritual Shap roamed off to fanciful lands. As he walked from the station, dreamy but wide awake, it suddenly struck him that the real Shap was not the one walking to Business in black and ugly clothes, but he who roamed along a jungle's edge near the \glos{ramparts}{defensive walls around a city or castle}{noun} of an old and Eastern city that rose up sheer from the sand, and against which the desert lapped with one eternal wave. He used to fancy the name of that city was Larkar. "After all, the fancy is as real as the body," he said with perfect logic. It was a dangerous theory.

For that other life that he led he realized, as in Business, the importance and value of method. He did not let his fancy roam too far until it perfectly knew its first surroundings. Particularly he avoided the jungle—he was not afraid to meet a tiger there (after all it was not real), but stranger things might crouch there. Slowly he built up Larkar: rampart by rampart, towers for archers, gateway of brass, and all. And then one day he argued, and quite rightly, that all the silk-clad people in its streets, their camels, their \glos{wares}{goods offered for sale}{noun} that came from Inkustahn, the city itself, were all the things of his will—and then he made himself King. He smiled after that when people did not raise their hats to him in the street, as he walked from the station to Business; but he was sufficiently practical to recognize that it was better not to talk of this to those that only knew him as Mr. Shap.

Now that he was King in the city of Larkar and in all the desert that lay to the East and North he sent his fancy to wander further afield. He took the regiments of his camel-guards and went jingling out of Larkar, with little silver bells under the camels' chins, and came to other cities far-off on the yellow sand, with clear white walls and towers, uplifting themselves in the sun. Through their gates he passed with his three silken regiments, the light-blue regiment of the camel-guards being upon his right and the green regiment riding at his left, the lilac regiment going on before. When he had gone through the streets of any city and observed the ways of its people, and had seen the way that the sunlight struck its towers, he would proclaim himself King there, and then ride on in fancy. So he passed from city to city and from land to land. Clear-sighted though Mr. Shap was, I think he overlooked the lust of \glos{aggrandizement}{the desire for greater power or status}{noun} to which kings have so often been victims; and so it was that when the first few cities had opened their gleaming gates and he saw peoples prostrate before his camel, and spearmen cheering along countless balconies, and priests come out to do him reverence, he that had never had even the lowliest authority in the familiar world became unwisely insatiate. He let his fancy ride at inordinate speed, he forsook method, scarce was he king of a land but he yearned to extend his borders; so he journeyed deeper and deeper into the wholly unknown. The concentration that he gave to this inordinate progress through countries of which history is ignorant and cities so fantastic in their \glos{bulwarks}{strong defensive walls or protections}{noun} that, though their inhabitants were human, yet the foe that they feared seemed something less or more; the amazement with which he beheld gates and towers unknown even to art, and furtive people thronging intricate ways to acclaim him as their sovereign—all these things began to affect his capacity for Business. He knew as well as any that his fancy could not rule these beautiful lands unless that other Shap, however unimportant, were well sheltered and fed: and shelter and food meant money, and money, Business. His was more like the mistake of some gambler with cunning schemes who overlooks human greed. One day his fancy, riding in the morning, came to a city gorgeous as the sunrise, in whose opalescent wall were gates of gold, so huge that a river poured between the bars, floating in, when the gates were opened, large \glos{galleons}{large old sailing ships}{noun} under sail. Thence there came dancing out a company with instruments, and made a melody all around the wall; that morning Mr. Shap, the bodily Shap in London, forgot the train to town.

Until a year ago he had never imagined at all; it is not to be wondered at that all these things now newly seen by his fancy should play tricks at first with the memory of even so sane a man. He gave up reading the papers altogether, he lost all interest in politics, he cared less and less for things that were going on around him. This unfortunate missing of the morning train even occurred again, and the firm spoke to him severely about it. But he had his consolation. Were not Arathrion and Argun Zeerith and all the level coasts of Oora his? And even as the firm found fault with him his fancy watched the \glos{yaks}{long-haired pack animals from Central Asia}{noun} on weary journeys, slow specks against the snow-fields, bringing tribute; and saw the green eyes of the mountain men who had looked at him strangely in the city of Nith when he had entered it by the desert door. Yet his logic did not forsake him; he knew well that his strange subjects did not exist, but he was prouder of having created them with his brain, than merely of ruling them only; thus in his pride he felt himself something more great than a king, he did not dare to think what! He went into the temple of the city of Zorra and stood some time there alone: all the priests kneeled to him when he came away.

He cared less and less for the things we care about, for the affairs of Shap, the business-man in London. He began to despise the man with a royal contempt.

One day when he sat in Sowla, the city of the Thuls, throned on one \glos{amethyst}{a purple precious stone}{noun}, he decided, and it was proclaimed on the moment by silver trumpets all along the land, that he would be crowned as king over all the lands of Wonder.

By that old temple where the Thuls were worshipped, year in, year out, for over a thousand years, they pitched pavilions in the open air. The trees that blew there threw out radiant scents unknown in any countries that know the map; the stars blazed fiercely for that famous occasion. A fountain hurled up, clattering, ceaselessly into the air armfuls on armfuls of diamonds. A deep hush waited for the golden trumpets, the holy coronation night was come. At the top of those old, worn steps, going down we know not whither, stood the king in the emerald-and-amethyst cloak, the ancient garb of the Thuls; beside him lay that Sphinx that for the last few weeks had advised him in his affairs.

Slowly, with music when the trumpets sounded, came up towards him from we know not where, one-hundred-and-twenty archbishops, twenty angels and two archangels, with that terrific crown, the \glos{diadem}{a crown or royal headpiece}{noun} of the Thuls. They knew as they came up to him that promotion awaited them all because of this night's work. Silent, majestic, the king awaited them.

The doctors downstairs were sitting over their supper, the \glos{warders}{attendants or guards in an institution}{noun} softly slipped from room to room, and when in that cosy dormitory of Hanwell they saw the king still standing erect and royal, his face resolute, they came up to him and addressed him:

"Go to bed," they said—"pretty bed." So he lay down and soon was fast asleep: the great day was over.\source{https://www.gutenberg.org/files/7477/7477-h/7477-h.htm}
\end{document}
